\documentclass[../cours_1re.tex]{subfiles}
\begin{document}
\setcounter{chapter}{6}
\chapter{Variations et limites de suites}
\label{ch:07}
\begin{objectif}
Étudier les variations d'une suite et distinguer convergence, divergence et oscillation. Le programme exige une approche intuitive des limites ; les définitions formelles ci-dessous sont un approfondissement.
\end{objectif}
\section{Comparer deux termes consécutifs}
Une suite est croissante si $u_{n+1}\ge u_n$ pour tout rang, décroissante si l'inégalité est inversée. Elle est majorée s'il existe un même réel $M$ tel que $u_n\le M$ pour tous les rangs ; minorée de même ; bornée si les deux propriétés sont réunies.
\begin{methode}
Étudier $u_{n+1}-u_n$. Pour des termes strictement positifs, comparer $u_{n+1}/u_n$ à $1$ est aussi possible. Une formule $u_n=f(n)$ permet d'utiliser les variations de $f$ sur un intervalle contenant les indices.
\end{methode}
\begin{exemple}
Pour $u_n=n/(n+1)$,
$u_{n+1}-u_n=1/((n+1)(n+2))>0$ et $0\le u_n<1$.
Pour $v_n=(-1)^n$, les valeurs restent entre $-1$ et $1$, mais alternent : être bornée ne signifie pas converger.
\end{exemple}
\section{Ce que signifie se rapprocher}
$1/n$ se rapproche de zéro : on peut rendre tous les termes à partir d'un certain rang aussi petits qu'on le souhaite.
$n^2$ dépasse durablement n'importe quelle borne fixée.
La suite $(-1)^n$ ne se rapproche d'aucun nombre ; elle ne tend pas non plus vers $+\infty$ ou $-\infty$.
\begin{definition}[Complément : limite finie]
La suite $(u_n)$ tend vers $\ell$ si, pour tout $\varepsilon>0$, il existe un rang $N$ tel que, pour tout $n\ge N$, $|u_n-\ell|<\varepsilon$.
On écrit $u_n\to\ell$.
\end{definition}
Le rang $N$ peut dépendre de la précision $\varepsilon$, mais ne dépend pas de $n$ après son choix.
\begin{demonstration}[Deux exemples]
Pour $u_n=1/n$, choisir $N>1/\varepsilon$ suffit.
Pour $u_n=n/(n+1)=1-1/(n+1)$, le même raisonnement prouve la limite $1$. L'identité donne l'erreur exacte, pas seulement une intuition graphique.
\end{demonstration}
\section{Unicité et suites extraites}
\begin{theoreme}[Complément]
Une suite convergente a une seule limite et elle est bornée. Toute suite extraite $(u_{n_k})$, avec $n_k$ strictement croissant, converge vers la même limite.
\end{theoreme}
\begin{demonstration}
Si $\ell\ne m$ étaient deux limites, prendre $\varepsilon=|\ell-m|/3$. Pour un rang assez grand,
$|\ell-m|\le|\ell-u_n|+|u_n-m|<2|\ell-m|/3$, contradiction.
Pour la borne, tous les termes assez avancés sont à distance inférieure à $1$ de $\ell$ ; les premiers termes sont en nombre fini, donc bornés.
Enfin $n_k\ge k+n_0$ finit par dépasser tout rang fixé : la condition de convergence s'applique aux termes extraits.
\end{demonstration}
Les sous-suites de rangs pairs et impairs de $(-1)^n$ ont les limites distinctes $1$ et $-1$ : la suite complète ne converge pas.
\section{Une preuve du comportement géométrique}
Pour $0<q<1$, écrire $1/q=1+t$ avec $t>0$. L'inégalité de Bernoulli donne
\[
0<q^n=\frac1{(1+t)^n}\le\frac1{1+nt}.
\]
Le membre de droite devient arbitrairement petit : $q^n\to0$.
Pour $-1<q<0$, $|q^n|=|q|^n\to0$, donc la limite est encore zéro.
Pour $q>1$, $q^n\ge1+n(q-1)$ dépasse toute borne ; pour $q\le-1$, les signes alternent avec amplitude constante ou croissante.
\begin{attention}
Une suite croissante peut converger : $n/(n+1)\to1$. Une suite non monotone peut converger : $(-1/2)^n\to0$. « Augmenter » et « tendre vers l'infini » ne sont pas synonymes.
\end{attention}
\section{Passer à une limite : justification}
\approfondissement Si $u_n\to u$ et $v_n\to v$, alors $u_n+v_n\to u+v$ :
$|(u_n+v_n)-(u+v)|\le|u_n-u|+|v_n-v|$, et chaque erreur peut être rendue inférieure à $\varepsilon/2$.
Pour le produit,
$u_nv_n-uv=u_n(v_n-v)+v(u_n-u)$ ; la suite $u_n$ est bornée et les deux erreurs tendent vers zéro.
Pour le quotient, il faut $v\ne0$ : pour $n$ assez grand, $|v_n|\ge|v|/2$, puis
$|1/v_n-1/v|=|v_n-v|/(|v_nv|)\to0$.
\begin{demonstration}[Encadrement]
Si $|u_n-\ell|\le a_n$ à partir d'un rang et $a_n\to0$ avec $a_n\ge0$, choisir un rang où $a_n<\varepsilon$ donne directement $u_n\to\ell$.
\end{demonstration}
\begin{finchapitre}
\section{Monotonie et complétude}
\begin{theoreme}[Complément]
Toute suite croissante et majorée converge.
\end{theoreme}
\begin{demonstration}
On utilise la propriété de borne supérieure des réels : l'ensemble non vide et majoré $\{u_n:n\in\N\}$ possède une borne supérieure $\ell$. Pour $\varepsilon>0$, $\ell-\varepsilon$ n'est pas un majorant ; il existe donc $N$ tel que $\ell-\varepsilon<u_N$. La croissance donne, pour $n\ge N$,
$\ell-\varepsilon<u_N\le u_n\le\ell$. C'est la convergence.
\end{demonstration}
La propriété de borne supérieure est un fondement admis de $\R$, pas une conséquence de quelques essais numériques.
Pour $u_{n+1}=\sqrt{2+u_n}$, $u_0=1$, la stabilité de $[1;2]$ et
$\sqrt{2+x}\ge x$ sur cet intervalle donnent croissance et convergence. En passant à la limite, $\ell^2=2+\ell$ ; les solutions sont $-1$ et $2$, mais $\ell\in[1;2]$ impose $\ell=2$.
Résoudre l'équation du point fixe seul n'aurait pas prouvé l'existence de la limite.


\section{Utilité et prolongements}
Les limites servent à certifier qu'un calcul itératif approche réellement une valeur. Le contrôle $|u_n-\ell|\le Cq^n$ fournit aussi un critère d'arrêt.

En licence, la complétude distingue $\R$ de $\Q$ et fonde les théorèmes d'existence. En master, on étudie des suites de fonctions ou d'opérateurs : il faut préciser la notion de distance et le mode de convergence. Le schéma « majorer une erreur par une quantité qui tend vers zéro » restera central.
\begin{retenir}
Une équation de limite ne prouve pas la convergence. Séparer observation, existence de la limite, identification de sa valeur et estimation de l'erreur.
\end{retenir}
\end{finchapitre}
\end{document}
